% GENERATED FILE. Edit canonical lessons, then run npm run generate:books.
% canonical-source-hash: fnv1a64:185654f2a3ad4e17
% canonical-lessons: GE-C271-melden, GE-C271-meldung, GE-C271-bericht, GE-C271-schlagzeile, GE-C271-artikel

\chapter{Die Meldung}
\label{ch:ge-die-meldung}
\hlchaptermodality{271}

\begin{chapteropening}
\textbf{By the end of this chapter:} \emph{I can talk about the news: to report, a news item, a report, a headline and an article.}

Complete the last lesson of chapter 271: I can talk about the news: to report, a news item, a report, a headline and an article.
\end{chapteropening}

\section[melden]{melden — to report, to announce; reflexive: to get in touch}

\label{lesson:GE-C271-melden}



\begin{quote}
Before the new one: say the German for to go, then the German for the sequence of events.
\end{quote}

\subsection*{You'll want to know: melden}
\textbf{melden} — \textquotedblleft{}to report, to announce; reflexive: to get in touch\textquotedblright{}.

\begin{quote}
Die Zeitung meldete einen großen Unfall.
\end{quote}

\begin{quote}
\emph{The newspaper reported a big accident.}
\end{quote}

Reflexive, it means to get in touch: \emph{Ich melde mich morgen}.

\begin{cousinweb}
\emph{melden} is inherited from Old High German \emph{meldōn}, which EtymDB traces to Proto-Germanic \emph{melþōną} \textquotedblleft{}to make known, to give away\textquotedblright{}.
\end{cousinweb}

\begin{grammarlens}[title={What you've built}]
One more word for telling, explaining and describing at B1.
\end{grammarlens}

\subsection*{Your turn}
\begin{itemize}
\raggedright
  \item \emph{Say it:} \emph{melden}
  \item \emph{Say it:} \emph{melden}, once more
  \item \emph{Say it:} a sentence of your own with \emph{melden}
  \item \emph{Recall:} say the German for to go, then the German for the sequence of events, then say \emph{melden} again
\end{itemize}

\begin{tcolorbox}[breakable,colback=teal!4,colframe=teal!35!black,title={Before you move on}]
What does melden mean? (\textquotedblleft{}To report, to announce; reflexive: to get in touch\textquotedblright{}.) Say it once more.
\end{tcolorbox}
\section[die Meldung]{die Meldung — the news item, the announcement}

\label{lesson:GE-C271-meldung}



\begin{quote}
Before the new one: say the German for the sequence of events, then the German for to report.
\end{quote}

\subsection*{You'll want to know: die Meldung}
\textbf{die Meldung} — \textquotedblleft{}the news item, the announcement\textquotedblright{}.

\begin{quote}
Die Meldung kam am Abend im Radio.
\end{quote}

\begin{quote}
\emph{The news item came on the radio in the evening.}
\end{quote}

Plural: \emph{die Meldungen}. Also an official notice to an office.

\begin{cousinweb}
\emph{Meldung} is built on \emph{melden} \textquotedblleft{}to report\textquotedblright{}, inherited from Old High German \emph{meldōn} (EtymDB), with the noun ending \emph{-ung}.
\end{cousinweb}

\begin{grammarlens}[title={What you've built}]
One more word for telling, explaining and describing at B1.
\end{grammarlens}

\subsection*{Your turn}
\begin{itemize}
\raggedright
  \item \emph{Say it:} \emph{die Meldung}
  \item \emph{Say it:} \emph{die Meldung}, once more
  \item \emph{Say it:} a sentence of your own with \emph{die Meldung}
  \item \emph{Recall:} say the German for the sequence of events, then the German for to report, then say \emph{die Meldung} again
\end{itemize}

\begin{tcolorbox}[breakable,colback=teal!4,colframe=teal!35!black,title={Before you move on}]
What does die Meldung mean? (\textquotedblleft{}The news item, the announcement\textquotedblright{}.) Say it once more.
\end{tcolorbox}
\section[der Bericht]{der Bericht — the report}

\label{lesson:GE-C271-bericht}



\begin{quote}
Before the new one: say the German for to report, then the German for the news item.
\end{quote}

\subsection*{You'll want to know: der Bericht}
\textbf{der Bericht} — \textquotedblleft{}the report\textquotedblright{}.

\begin{quote}
Der Bericht war sehr lang.
\end{quote}

\begin{quote}
\emph{The report was very long.}
\end{quote}

Plural: \emph{die Berichte}. Built on \emph{berichten}, already known.

\begin{cousinweb}
\emph{Bericht} comes from Middle High German \emph{beriht}, built on \emph{berihten} \textquotedblleft{}to put in order, to instruct\textquotedblright{} (EtymDB), a relative of \emph{richten} and \emph{recht}.
\end{cousinweb}

\begin{grammarlens}[title={What you've built}]
One more word for telling, explaining and describing at B1.
\end{grammarlens}

\subsection*{Your turn}
\begin{itemize}
\raggedright
  \item \emph{Say it:} \emph{der Bericht}
  \item \emph{Say it:} \emph{der Bericht}, once more
  \item \emph{Say it:} a sentence of your own with \emph{der Bericht}
  \item \emph{Recall:} say the German for to report, then the German for the news item, then say \emph{der Bericht} again
\end{itemize}

\begin{tcolorbox}[breakable,colback=teal!4,colframe=teal!35!black,title={Before you move on}]
What does der Bericht mean? (\textquotedblleft{}The report\textquotedblright{}.) Say it once more.
\end{tcolorbox}
\section[die Schlagzeile]{die Schlagzeile — the headline}

\label{lesson:GE-C271-schlagzeile}



\begin{quote}
Before the new one: say the German for the news item, then the German for the report.
\end{quote}

\subsection*{You'll want to know: die Schlagzeile}
\textbf{die Schlagzeile} — \textquotedblleft{}the headline\textquotedblright{}.

\begin{quote}
Die Schlagzeile war groß und rot.
\end{quote}

\begin{quote}
\emph{The headline was big and red.}
\end{quote}

Plural: \emph{die Schlagzeilen}. \emph{Schlagzeilen machen} — to hit the headlines.

\begin{cousinweb}
\emph{Schlagzeile} joins \emph{Schlag} \textquotedblleft{}a blow\textquotedblright{}, from \emph{schlagen} \textquotedblleft{}to strike\textquotedblright{}, and \emph{Zeile} \textquotedblleft{}a line\textquotedblright{}; both are inherited Germanic words. EtymDB lists no parent for the compound.
\end{cousinweb}

\begin{grammarlens}[title={What you've built}]
One more word for telling, explaining and describing at B1.
\end{grammarlens}

\subsection*{Your turn}
\begin{itemize}
\raggedright
  \item \emph{Say it:} \emph{die Schlagzeile}
  \item \emph{Say it:} \emph{die Schlagzeile}, once more
  \item \emph{Say it:} a sentence of your own with \emph{die Schlagzeile}
  \item \emph{Recall:} say the German for the news item, then the German for the report, then say \emph{die Schlagzeile} again
\end{itemize}

\begin{tcolorbox}[breakable,colback=teal!4,colframe=teal!35!black,title={Before you move on}]
What does die Schlagzeile mean? (\textquotedblleft{}The headline\textquotedblright{}.) Say it once more.
\end{tcolorbox}
\section[der Artikel]{der Artikel — the article}

\label{lesson:GE-C271-artikel}



\begin{quote}
Before the new one: say the German for the report, then the German for the headline.
\end{quote}

\subsection*{You'll want to know: der Artikel}
\textbf{der Artikel} — \textquotedblleft{}the article\textquotedblright{}.

\begin{quote}
Ich habe einen interessanten Artikel über Bienen gelesen.
\end{quote}

\begin{quote}
\emph{I read an interesting article about bees.}
\end{quote}

Plural: \emph{die Artikel}. Also the grammar word: \emph{der}, \emph{die} and \emph{das} are \emph{Artikel}.

\begin{cousinweb}
\emph{Artikel} is a learned borrowing from Latin \emph{articulus} \textquotedblleft{}a joint, a part, the article in grammar\textquotedblright{} (EtymDB), a diminutive of \emph{artus} \textquotedblleft{}joint\textquotedblright{} (Lewis \& Short).
\end{cousinweb}

\begin{grammarlens}[title={What you've built}]
One more word for telling, explaining and describing at B1.
\end{grammarlens}

\subsection*{Your turn}
\begin{itemize}
\raggedright
  \item \emph{Say it:} \emph{der Artikel}
  \item \emph{Say it:} \emph{der Artikel}, once more
  \item \emph{Say it:} a sentence of your own with \emph{der Artikel}
  \item \emph{Recall:} say the German for the report, then the German for the headline, then say \emph{der Artikel} again
\end{itemize}

\begin{tcolorbox}[breakable,colback=teal!4,colframe=teal!35!black,title={Before you move on}]
What does der Artikel mean? (\textquotedblleft{}The article\textquotedblright{}.) Say it once more.
\end{tcolorbox}
